\section{Discussion}
\label{s:discussion}

Placing parallel trends on the log reproduction number with imported infection pressure, and
eliminating unrestricted unit effects by forward quasi-differencing, yields a difference-in-differences
design that accommodates self-exciting transmission, incidence-driven policy feedback and lagged
spatial contagion while retaining closed-form week profiling and the ordinary cluster-robust
sandwich. Across 25 pre-specified simulation scenarios spanning feedback adoption, spillover,
overdispersion and small counts, QD-GMM recovers direct and spillover effects with negligible bias
and nominal coverage where count, log-count, Poisson fixed-effects and two-step reproduction-number
comparators fail, and directly identifies the policy-level total $\tau+\kappa$ when interventions
are enacted above the reporting unit.

The German county-level varicella application demonstrates this framework in its intended operating
regime. Varicella's 14-day generation interval satisfies weekly aggregation ($w_0=2.6\times 10^{-6}$),
and staggered six-week summer school holidays across 16 federal states provide recurring,
non-absorbing exposure variation. QD-GMM identifies a unique interior contagion coupling
($\hat\phi=0.325$) and a 21.2\% reduction in transmission ($\hat\tau+\hat\kappa=-0.238$, 95\%
$t_{15}$ CI $-0.362$ to $-0.114$), matching European varicella benchmarks \citep{jackson2014varicella}
and corresponding to 2{,}440 cases averted in treated county-weeks and 2{,}995 network-wide. The
finding holds across score wild bootstrap-$t$, nonlinear pairs cluster bootstrap-$t$,
delete-one-state jackknife, NUTS-2 spatial aggregation and the identification-robust Stock--Wright
Bonferroni projection ($-0.492$ to $-0.081$), while event-time estimates exhibit flat pre-trends and
a null onset week ($-0.008$) followed by a post-incubation drop. Equally instructive, standard
comparators reproduce Theorem~\ref{thm:bias} on real data: two-step EpiEstim log-$R$ regression
flips to the wrong sign ($+0.064$) under 58.8\% county zeros before reverting to $-0.102$ at NUTS-2
aggregation (2.9\% zeros), and $\log(Y+1)$ TWFE is attenuated to $-0.102$.

Three practical conditions define the scope of weekly QD-GMM. First, global identification relies on
a unique interior root for $\phi$: at the zero-contagion boundary ($\phi_0=0$), the population
moments can admit a second $\bar\phi$-dependent root. Inspecting the profiled $\phi$-equation over
$[0,3]$ and reporting Stock--Wright identification-robust sets alongside M0-G and M0-GI protects
against spurious roots without assuming strong identification of $\phi$ \citep{stock2000weak}.
Second, weekly aggregation requires that within-week transmission be small: the matched stress
experiment ($w_0=0.3$; Section~\ref{ss:sameweek}) shows that substantial same-week renewal biases
weekly policy totals by roughly one-sixth while leaving $J$ unpowered. Weekly QD-GMM is therefore
well suited to multi-week generation-interval pathogens such as varicella, measles or mumps, whereas
fast-generation pathogens such as SARS-CoV-2 require sub-weekly estimating equations that also model
day-of-week reporting and notification delays \citep{lison2024nowcast}. Third, Theorem~\ref{thm:asy}
is a large-$G$ result: with $G=16$ unequal state clusters (1 to 96 counties), calibrated simulations
show mild under-coverage of $t_{15}$ for $\tau+\kappa$ (0.892--0.902), so corroborating sandwich
intervals with nonlinear pairs cluster bootstrap, delete-one-cluster jackknife and identification-robust
projections is important whenever clusters are moderate in number or heterogeneous in size
\citep{mackinnon2017wild,ferman2019inference,canay2021wild}. Similarly, in sparse-count settings
where mean incidence is comparable to the background term $\omega_i$, sensitivity to $\omega_i$ and
score-form over-identification diagnostics ($J_{LM,\mathrm{inv}}$, $J_{LM,\mathrm{oi}}$) should be
inspected to separate short-run reporting noise from structural misspecification.

Framing parallel trends on $\log R$ \citep{roth2023parallel} separates local from imported infection
pressure \citep{elenev2025staggered} and complements unconfoundedness approaches
\citep{callaway2023pandemic}. By turning routine county-week surveillance counts into
transmission-scale policy estimates with standard clustered inference and transparent falsification
diagnostics, QD-GMM provides a practical causal tool for infectious-disease epidemiology. Natural
extensions include sequential efficiency bounds \citep{chamberlain1992sequential,chamberlain2022feedback},
likelihood-based feedback models \citep{bonhomme2025feedback}, sub-weekly estimating equations,
reporting-independent signals such as wastewater surveillance \citep{goldstein2024wastewater}, and
joint estimation of generation intervals and reporting processes \citep{goldwasser2026rt}.
